\movedhead{A format that was rejected}
\label{sec:rejected}\label{app:rejected}
A forced choice between a winner and a loser is the canonical
operationalisation of the disposition effect. Asked to raise cash by selling
either position A or position B, identical but for cost basis, the model
chose A in 24 of 24 trials irrespective of which position was the winner,
which presents as a large position preference. Offered a continuous
allocation over the same two positions, all 72 responses answered exactly
50.0 with zero variance. The apparent preference was an artifact of forcing a
binary choice on a model that had correctly identified the options as
equivalent, and is reported here as a warning about binary forced choice
items rather than as a finding about order effects. No forced binary choice
appears anywhere in the study.

The two asset format is not a straw man: it is what the largest audit in this
area uses. \citet{malberg2025comprehensive} present two owned assets, one up
and one down, ask which the model would rather sell, and swap which asset has
gained between the control and treatment arms, scoring a disposition effect
only when the preference for selling the winner survives the swap. That swap
is precisely what protects their design from the artifact above, since a
model that names the same position in both arms scores near zero rather than
as biased, and their graded 7-point response avoids the hard binary we
rejected. What the format cannot do, in either version, is attach a cost
basis to a single position and ask whether to act on it. A relative
preference between two assets is not a realisation decision, and it is the
realisation decision that the human measurements of
\citet{shefrin1985disposition} and \citet{odean1998losses} are about.
